\section{Background}
\label{sec:background}

\subsection{Where the metrics come from}

Vision-based accident anticipation was posed in its modern form by Chan
\textit{et al.}~\cite{chan2016dad}, who introduced the Dashcam Accident Dataset
together with the measure the field still reports: the interval between the
first frame at which a risk score crosses a threshold and the annotated onset
of the collision. Bao \textit{et al.}~\cite{bao2020ccd} contributed the Car
Crash Dataset and an uncertainty-based formulation, and Karim \textit{et
al.}~\cite{karim2022dsta} gave the reporting convention most subsequent work
adopts: average precision, mean time-to-accident over thresholds, and
time-to-accident at 80\% recall.

It is worth stating plainly what that convention gets right, because this paper
is about what happens when it is set aside. It reports a discrimination metric
alongside every timing metric, and it reads timing \emph{at a fixed operating
point} of the discrimination metric. A method therefore cannot buy
time-to-accident by alarming indiscriminately: the recall constraint holds it
to a fixed false-alarm behaviour before its timing is counted at all. The
repair proposed in Section~\ref{sec:protocol} is in large part a return to that
discipline.

Architectures built on this foundation run from supervised convolutional and
temporal models~\cite{chan2016dad, bao2020ccd, karim2022dsta} through
self-supervised formulations that reduce the dependence on scarce collision
labels~\cite{yao2019unsupervised, liu2026sctnet, zhong2025early}, to
training-free pipelines pairing a vision--language
model~\cite{radford2021clip} with a motion
signal~\cite{thakur2026modular, saha2026zeroshot, singh2026metadata} --- though
those target detection or classification and none reports a timing metric.

Two lines of work deserve particular attention, because they anticipate part of
this paper's subject. Pjetri \textit{et al.}~\cite{pjetri2025nondecreasing}
train with an objective requiring the predicted danger score to be
non-decreasing as the collision approaches, and Zou \textit{et
al.}~\cite{zou2026riskprop} introduce an adaptive monotonic constraint loss with
the same intent. Both are well-motivated responses to score oscillation and we
criticise neither. But both illustrate the hazard examined here: a risk score
constrained not to fall is, by construction, closer to the maximiser of a
threshold-crossing timing metric than a score free to fall. The metric cannot
separate the benefit of the constraint from the artefact of it, and neither
paper reports a control that would.

\subsection{From a convention to a single number}

Competitions need something the convention of~\cite{karim2022dsta} does not
supply: a total order over submissions. The natural way to obtain one is to sum
the metrics, and that is what the benchmark studied here
does~\cite{autopilot2026}, adding two threshold-crossing timing terms to
average precision and area under the ROC curve.

Summing is where the difficulty enters, and it is a difficulty of type rather
than of tuning. Average precision and area under the curve are bounded on the
unit interval; a time-to-accident measured in frames is bounded only by the
length of the clip. Adding them makes the sum depend on how the corpus was
windowed and lets the unbounded term decide the ranking.
Section~\ref{sec:why} develops the argument; Sections~\ref{sec:leaderboard}
and~\ref{sec:nexar} measure its magnitude.

\subsection{The missing control}

The check that exposes such a problem is a control that does not read the
input, and its value has been established repeatedly elsewhere: dataset
signatures strong enough to identify a photograph's
source~\cite{torralba2011bias}; models given only half the input scoring far
above chance on standard inference
corpora~\cite{gururangan2018artifacts, poliak2018hypothesis}, after which
partial-input baselines became standard practice; and tuned simple baselines
matching most recently published neural
recommenders~\cite{ferraridacrema2019progress}. In each case the control was
cheap, and in each case nobody had run it.

Accident anticipation has every precondition those episodes had. It reports a
timing metric whose maximiser is trivial, it is moving toward composite
single-number scores, and its newest benchmarks withhold labels --- so an
entrant cannot compute the discrimination half at all and must develop against
a locally computable proxy, the only candidate being the shape of one's own
risk curve, which a constant maximises.

The pathology has not gone unnoticed, but it has been noticed as an anomaly
rather than isolated. Zhao \textit{et al.}~\cite{zhao2025top} observe that the
standard calculation credits a false alarm in a safe scene with the interval to
some later, causally unrelated collision. Goldshmidt \textit{et
al.}~\cite{goldshmidt2025badas} report baselines making physically implausible
predictions nine to ten seconds ahead. Most directly, Caselli \textit{et
al.}~\cite{caselli2026flara} publish a table in which two established methods
record areas under the curve of $45.5$ and $41.0$ --- below chance --- beside
mean times-to-accident of $10.216$ and $10.263$ seconds. A method that cannot
separate accident clips from ordinary ones is there credited with ten seconds
of anticipation: this paper's argument, already legible in someone else's
results table. None of these accounts isolates the effect, and none reports
what a submission that reads nothing would score.

The nearest methodological precedent lies in an adjacent task. Korkut
\textit{et al.}~\cite{korkut2026blind} audit four traffic video
question-answering benchmarks by running vision--language models blind, and find
that removing the video consistently improves accuracy on one of them. That is
a control of exactly the kind advocated here, applied to multiple-choice
accuracy rather than to per-frame risk. We are not aware of an equivalent
control reported for the timing and discrimination metrics used in
anticipation, which is the gap this paper fills.

We state the scope of that observation precisely: it covers the works cited
here whose results tables we inspected directly. Four we consulted are behind
paywalls we could not open, and we exclude them rather than assume. It is not a
systematic audit of the field and we do not present it as one.

\subsection{What is and is not new here}

Monotonicity constraints on risk scores are prior
work~\cite{pjetri2025nondecreasing, zou2026riskprop} and this paper claims
none; so are training-free pipelines pairing a vision--language model with a
motion signal~\cite{thakur2026modular}. Nor is the general idea that a
benchmark can be satisfied by a degenerate submission new, and the suspicion
that time-to-accident is inflatable has been voiced
before~\cite{zhao2025top, goldshmidt2025badas, caselli2026flara}, as has a blind
control on an adjacent traffic-video task~\cite{korkut2026blind}. What is new is
the measurement of how far it goes on this class of metric, a probe methodology
that decomposes such a score from outside the competition holding the labels,
and an evaluation protocol that has been tested against adversarial controls,
found insufficient, and repaired.
